\scanpage{37}
\mspage{18}
$F(\Delta)/\{1,\underline a_\Delta\}$ (cf (36)), qui est donc l'\unsure{ens|une} des quatre bifaces de $\Delta$. Donc c'est l'octaèdre orienté canoniquement associé à l'octaèdre $\Delta$, savoir

\[
\left\{\begin{aligned}
\Delta^+ &= \deleted{\Delta\wedge_{\unsure{\{1,\underline a_\Delta\}}}}
 \mathrm{or}(\Delta)\wedge_{\{1,\underline a_\Delta\}}\Delta\\
\Delta &= \mathrm{or}(\Delta)\wedge_{\{1,\underline a_\Delta\}}\Delta^+
\end{aligned}\right. \tag*{(45)}
\]

de sorte qu'on a \inserted{bien}, comme il se doit

\[
\left\{\begin{aligned}
\mathrm{or}(\Delta^+) &\simeq \{\pm1\}=\mathbb 1_{\mathbb F_2}\\
\underline\omega_{\Delta^+} &\simeq \underline\omega_\Delta\simeq\underline\omega_J=\underline\omega
\quad [\text{car }\underline a_\Delta\text{ opère trivialement sur }\underline\omega_\Delta]\\
C_{\Delta^+} &\simeq C_\Delta\wedge\mu \simeq \underline\omega\qquad \ldots
\end{aligned}\right. \tag*{(45)}
\]

$[\![$ Il importe maintenant de bien comprendre la structure d'extension

\diagram{\begin{tikzcd}[ampersand replacement=\&,column sep=small,row sep=0pt]
{(47)}\&1\arrow[r] \& \mathfrak z\arrow[r] \& \mathrm{Gl}(M)\arrow[r] \& \mathrm{Gl}(M)'\arrow[r,phantom,"{\unsure{\text{---}|\to}}" description] \& 1\\
\&\& \rotatebox[origin=c]{90}{$\scriptstyle\simeq$} \& \& \rotatebox[origin=c]{90}{$\scriptstyle\simeq$} \&\\
\&\& \{\pm1\} \& \& \mathrm{Aut}_{\mathrm{oct}}(\Delta_M)' \&\\
\&\& \& \& \simeq\mathrm{Aut}_{\mathrm{oct}}^+(\Delta_M)\times\{1,\underline a_\Delta\}_{\unsure{M|\mu}} \&
\end{tikzcd}}

et en particulier

\diagram{\begin{tikzcd}[ampersand replacement=\&,column sep=small,row sep=0pt]
{(48)}\&1\arrow[r] \& \mathfrak z_0\arrow[r] \& \mathrm{Gl}(2,\mathbb Z/4\mathbb Z)\arrow[r] \& \mathrm{Gl}(2,\mathbb Z/4\mathbb Z)'\arrow[r,phantom,"{\unsure{\text{---}|\to}}" description] \& 1\\
\&\& \shortparallel \& \& \rotatebox[origin=c]{90}{$\scriptstyle\simeq$} \&\\
\&\& \{\pm1\} \& \& \mathrm{Aut}'_{\mathrm{oct}}\bigl(\Delta_{\underbrace{(\mathbb Z/4\mathbb Z)^2}_{\substack{\text{\textquotedblleft octaèdre}\\\text{\unsure{standart|standard}\textquotedblright}}}}\bigr) \&
\end{tikzcd}}

puisque c'est par là qu'on comprend quelle est la donnée supplémentaire qu'il faut pour déduire, d'un octaèdre non orienté $\Delta$, le $(\mathbb Z/4\mathbb Z)$-module libre de rang $2$ $M$ dont $\Delta$ proviendrait -- question qui revient à celle de \textquotedblleft remonter\textquotedblright\ un torseur sous le groupe $\mathrm{Gl}(2,\mathbb Z/4\mathbb Z)'$ en un torseur sous le groupe $\mathrm{Gl}(2,\mathbb Z/4\mathbb Z)$ lui-même.
% archive sheet number: 36
